\documentclass[10pt,a4paper]{amsart}
\usepackage[margin=19mm]{geometry}
\usepackage{amsmath,amssymb,mathtools}
\usepackage[scr=rsfs,cal=euler]{mathalfa}
\usepackage{xfrac}
\usepackage[unicode,hidelinks,bookmarks=false]{hyperref}
\newcommand{\ie}{i.e.\ }
\newcommand{\cf}{cf.\ }
\newcommand{\resp}{respectively\ }
\newcommand{\Subsection}{\subsection}
\providecommand{\nobreakdash}{\nobreak\hbox{-}}
\setlength{\emergencystretch}{2em}
\multlinegap0pt
\usepackage{bm}
\usepackage{bbm}
\usepackage{dsfont}
\usepackage{xspace}
\usepackage{cancel}
\usepackage{mathtools}
\usepackage{amsthm}
\makeatletter
\@ifundefined{lem}{\newtheorem{lem}{Lemma}}{}
\makeatother
\title[On the edge densities of normal, convex mosaics]{On the edge densities of normal, convex mosaics}
\author{M\'at\'e Kadlicsk\'o, Zsolt L\'angi, Shanxiang Lyu}
\date{Source: arXiv:2312.08050v3 (CC BY 4.0)}
\begin{document}
\maketitle

\noindent\emph{Source and licence.} On the edge densities of normal, convex mosaics. Version arXiv:2312.08050v3 (CC BY 4.0), distributed under the Creative Commons Attribution 4.0 International licence, as recorded in the arXiv licence metadata for this exact version and shown on the abstract page https://arxiv.org/abs/2312.08050v3. Full rights notice: licenses/arxiv-2312.08050-CC-BY-4.0.txt.\par\smallskip

\noindent\textit{Editorial scope.} Counted unit: the Lem whose source label is \texttt{lem:computations2} in the article (source0.tex lines 418--427, source label \texttt{lem:computations2}), delivered together with the article's own complete original proof of it (source0.tex lines 429--451, one \texttt{proof} environment of 23 lines). No mathematics is written, repaired or re-derived: statement and proof are the article's own text line for line. Required context retained uncounted: none. Every definition, display and label of those lines is retained unchanged. Omitted from this package and not redistributed: 1--417, 428--428, 452--769. Nothing inside a retained excerpt is omitted or altered: every retained body is delivered exactly as the article's lines are written, so no omission receipt of any kind accompanies this package. Citations: the retained excerpts carry no citation command at all, so no bibliography entry of the article is transcribed and no reference is redistributed. Reference adaptation: none was needed and none was made; every reference inside the delivered unit resolves inside it, so no unresolved reference or citation is produced. Printed slips kept literal: no printed word, symbol, hypothesis or number of the counted statement or of its proof was altered. Layout and class: the article's own class is \texttt{amsart} with packages including latexsym, amssymb, amsmath, amsopn, graphicx, color, epsfig, tikz, subcaption, array, hyperref, mathtools; the wrapper is the lane's pinned \texttt{amsart} editorial wrapper at 10pt on A4, which restates the theorem-like environments the retained text needs and adds the article's own macro definitions that the retained text needs (none), with extra wrapper packages (bm, bbm, dsfont, xspace, cancel, mathtools). The delivered numbering is the unit's own. Assets: no retained span contains a figure environment, a graphics-inclusion command, a PDF-page inclusion, a table or a TikZ picture; no image file is copied into this package, the package declares no asset dependency and no image file is redistributed with it. Licence: version arXiv:2312.08050v3 of this article is distributed under the Creative Commons Attribution 4.0 International licence, as recorded in the arXiv licence metadata for this exact version and shown on the abstract page https://arxiv.org/abs/2312.08050v3. A full rights notice is delivered at licenses/arxiv-2312.08050-CC-BY-4.0.txt. Characters: every retained excerpt is pure ASCII, so the delivered engine, XeLaTeX with its default Latin Modern text font, reports no missing character.
\begin{lem}\label{lem:computations2}
Let $C > 0$ and $\lambda \geq 1$. Then for each $(\tau_{12},\tau_{13},\tau_{14},\tau_{23},\tau_{24},0)\in \mathcal{X}_C$ we have 
\begin{equation}\label{ineq:g}
 f(\lambda \tau_{12},\tau_{13},\tau_{14},\tau_{23},\tau_{24},0) \leq \frac{16C^3 \lambda^3}{27(4\lambda - 1)^2},
\end{equation}
with equality if and only if
\[
\tau_{12} = \frac{4 \lambda - 3}{2\lambda}\tau_{13} = \ldots = \frac{4 \lambda - 3}{2\lambda} \tau_{24} = \frac{2 \lambda}{12 \lambda - 3}C.
\]
\end{lem}

\begin{proof}
Without loss of generality, we assume that $C=1$ and for simplicity, in the proof we denote the operator $\frac{\partial}{\partial \tau_{ij}}$ by $\partial_{ij}$.
Given that $\mathcal{X}_C$ is compact, the function $f$ attains its maximum on it at some point $(\tau^*,0) \in \mathcal{X}_C$ with $\tau^*=(\tau^*_{12}, \ldots, \tau_{24}^*)$.

\emph{Case 1}: all coordinates of $\tau^*$ are positive.\\
Then, by the Lagrange multiplier method, the gradient of $f$ at $\tau^*$ is parallel to the gradient of the function $\sum_{i < j, ij \neq 34} \tau_{ij}-1$ defining the condition, thus, all five partial derivatives of $f$ at $\tau^*$ are equal. 
From the equations $(\partial_{13} f)(\tau^*) = (\partial_{14} f)(\tau^*)$ and $(\partial_{23} f)(\tau^*) = (\partial_{24} f)(\tau^*)$ and by the assumption that all coordinates of $\tau^*$ are positive, we obtain that $\tau_{13}^* = \tau_{14}^*$ and $\tau_{23}^* = \tau_{24}^*$. Eliminating $\tau_{14}^*$ and $\tau_{24}^*$ from the partial derivatives, the equation $(\partial_{13} f)(\tau^*) = (\partial_{23} f)(\tau^*)$, and the fact that all $\tau_{ij}^*$ as well as $\lambda$ are positive, yields that $\tau_{13}^*=\tau_{23}^*$. Thus, we need to find the maximum of $f(\lambda \tau_{12},\tau_{13},\ldots,\tau_{13},0)$ under the condition that $\tau_{12} + 4 \tau_{13} = 1$. From this, by simple calculus methods, we obtain that $f(\lambda \tau_{12},\tau_{13},\ldots,\tau_{13},0) \leq \frac{16 \lambda^3}{27(4\lambda - 1)^2}$, with equality if and only if $\tau_{12} = \frac{4 \lambda - 3}{2\lambda}\tau_{13} = \frac{2 \lambda}{12 \lambda - 3}$.

\emph{Case 2}: $\tau_{12}^* = 0$ and all other $\tau_{ij}^*$ are positive.\\
Note that $f(0,\tau_{13},\ldots,\tau_{24},0) = (\tau_{13}+\tau_{24})\tau_{14}\tau_{23}+(\tau_{14}+\tau_{23})\tau_{13}\tau_{24}$. In this case, by the same argument as in Case 1, we obtain that $\tau_{13}^*=\ldots=\tau_{24}^*=\frac{1}{4}$, and $f(\tau^*)=\frac{1}{16}$. 

\emph{Case 3}: $\tau_{ij}^* = 0$ for some $ij \neq 12$, and all other $\tau_{st}^*$ are positive.\\
Without loss of generality, we may assume that $\tau_{24}^* = 0$. By a method similar to the one in Case 1 and using the condition that $\lambda \geq 1$, we obtain that $\tau_{12}^* = \frac{2(2\lambda-1)}{3(4\lambda-1)}$, $\tau_{13}^* = \tau_{14}^* = \frac{2\lambda}{3(4\lambda-1)}$, and $\tau_{23}^* = 0$, and $f(\tau^*)=\frac{4\lambda^2}{27(4\lambda-1)}$.

\emph{Case 4}, at least two $\tau_{ij}^*$ are zero.\\
We may assume that exactly two $\tau_{ij}^*$ are zero, and that there is no $s \in \{1,2,3,4\}$ such that the indices of all nonzero $\tau_{ij}^*$s contain $s$, as otherwise $f(\tau^*,0)=0$. Under this condition, depending on which $\tau_{ij}^*$ are not zero, we have that $f(\tau^*)=\frac{\lambda}{27}$ or $f(\tau^*)=\frac{1}{27}$, with equality if and only if all nonzero $\tau_{ij}^*$ are equal to $\frac{1}{3}$.

As a summary, we need to find
\[
\max \left \{ \frac{16 \lambda^3}{27(4\lambda - 1)^2}, \frac{1}{16}, \frac{4\lambda^2}{27(4\lambda-1)}, \frac{\lambda}{27}, \frac{1}{27} \right\}.
\]
Then an elementary computation shows that under the condition $\lambda \geq 1$, the maximum is $\frac{16 \lambda^3}{27(4\lambda - 1)^2}$, and the assertion follows.
\end{proof}

\end{document}
